
% 11. The Smudge on a Clean Windowpane: Biological Asymmetry.tex

\section{깨끗한 유리창의 얼룩, 생물학적 비대칭성}
우리는 오랫동안 저출산을 경제와 정책의 문제로만 취급해 왔다. 집값이 비싸고, 교육비가 많이 들며, 고용이 불안정하기 때문에 청년들이 아이를 낳지 못한다고 굳게 믿는다. 그래서 세계 각국의 정부는 천문학적인 보조금을 쏟아붓고 제도를 고치며 가족을 지원하지만, 고도로 발전한 선진국들일수록 출산율의 추락은 멈추지 않는다.
그러나 디스턴스(DISTANCE)의 서늘한 열역학적 렌즈를 들이대면, 저출산은 단순한 비용 문제나 사회적 이기심의 결과가 아님이 폭로된다. 그것은 완벽한 '평활화(Equalization)'를 향해 질주하는 문명의 궤적과, 결코 깎아낼 수 없는 인간의 '생물학적 번식 구조'가 동일한 트랙 위에서 정면충돌하며 빚어내는 거대한 물리학적 파열음이다. 

인류 문명은 오랜 역사에 걸쳐 신분, 권력, 그리고 남녀의 역할이라는 모든 인위적 '비대칭성'들을 매끄럽게 깎아내며 가장 평등하고 개방적인 사회를 완성해 왔다. 육체의 힘이 중요했던 시대에는 남녀가 서로 다른 트랙에서 살아가며 분리된 경쟁을 치렀지만, 산업혁명과 정보혁명이 물리적 장벽을 해체하면서 두 성별은 점차 완벽히 동일한 무대 위로 융합되었다. 이제 남녀는 같은 학교에서 공부하고, 같은 시험을 치르며, 같은 직장의 승진 체계 안에서 동일한 성공을 목표로 맹렬히 모멘텀을 부딪치며 경쟁한다.

이처럼 모든 것이 공정해진 극한의 대칭적 무대 위에서, 전대미문의 우주적 역설이 발생한다. 디스턴스 에세이는 이를 '깨끗한 유리창의 얼룩'이라는 뼈아픈 비유로 관통한다. 더러운 유리창에 묻은 수많은 얼룩들은 눈에 잘 띄지 않는다. 그러나 문명의 평활화 엔진이 그 유리창을 완벽하게 투명하고 깨끗하게 닦아내자, 그 위에 남아 있는 단 하나의 작은 얼룩은 누구에게나 섬뜩할 정도로 선명하고 거대하게 발견된다. 과거의 사회는 정치, 경제, 교육 등 곳곳에 불평등과 비대칭성이라는 얼룩이 가득했기에 특정한 차이가 부각되지 않았다. 그러나 모든 인위적 장벽이 제거된 오늘날의 텅 빈 무대 위에는 오직 단 하나의 얼룩만이 홀로 거대하게 남게 되었다.

그것은 바로, 오직 여성의 몸을 통해서만 이루어지는 '임신과 출산'이라는 생물학적 비대칭성이다.
이것은 어떠한 고결한 정치 제도로도 깎아낼 수 없고, 어떠한 경제 정책으로도 완전히 소거할 수 없는 가장 차가운 물리적 팩트다. 남녀가 완벽히 동일한 경기장에서 동일한 규칙표를 들고 뛰는 이 공정 경쟁의 투기장 속에서, 특정 성별에게만 부과되는 출산의 생물학적 부담은 더 이상 과거와 같은 숭고한 역할이나 희생이 아니다. 그것은 동일한 경쟁을 치러야 하는 주체에게 가해지는 거대한 '기회비용'이자, 완벽한 대칭의 무대 위에서 가장 잔혹하게 두드러지는 '선천적 불공정'으로 감각되기 시작한 것이다.

이 서늘한 역학적 진실은 최근 전 세계 선진국에서 공통적으로 관찰되는 '젊은 여성들의 급격한 진보화 현상'마저 완벽하게 해체해 버린다. 한국을 비롯해 북미, 서유럽 등에서 극명하게 나타나는 젊은 세대의 정치적 성향 차이(젠더 양극화)는 단순한 인터넷의 선동이나 혐오가 아니다. 사회적 평등과 공정성사회적 평등과 공정성이 확대되고 대칭화된 경쟁이 극대화될수록, 결코 깎아낼 수 없는 이 마지막 비대칭성(출산)은 더욱 날카로운 기회비용이자 거대한 '불공정의 얼룩'으로 감각될 수밖에 없다는 맹렬한 역학적 자각의 발현이다.

과거의 사회는 교육, 정치, 직업, 경제 등 모든 곳에 불평등이라는 얼룩이 가득 묻어 있는 지저분한 유리창이었다. 그래서 특정한 비대칭성(출산의 부담) 하나가 특별히 부각되지 않았다. 그러나 문명의 맹렬한 평활화 엔진이 그 유리창을 완벽하게 투명하고 매끄럽게 닦아내 버리자, 그 텅 빈 무대 위에 홀로 남겨진 번식의 생물학적 비대칭성은 이제 견딜 수 없는 거대한 모순으로 폭발하기 시작한 것이다.

동일한 교실에서 공부하고, 동일한 잣대로 평가받으며, 동일한 트랙 위에서 맹렬하게 성공의 모멘텀을 다투도록 설계된 이 완벽한 대칭성의 투기장 안에서, 오직 여성만이 짊어져야 하는 임신과 출산이라는 생물학적 조건은 더 이상 고결한 희생이나 자연의 섭리가 아니다. 그것은 공정 경쟁이라는 문명의 가장 숭고한 규칙을 정면으로 위반하는 치명적인 '선천적 기회비용'으로 인식될 수밖에 없는 것이다.

이 지점에서 인류가 그토록 찬미해 온 '공정성'의 서늘한 우주적 역설이 장엄하게 완성된다. 저출산은 단순히 청년들이 가난하거나 이기적이기 때문에 발생하는 경제적 파업이 아니다. 그것은 인류가 수천 년간 피 흘려 닦아온 '완벽한 공정성(평활화)'이라는 문명의 맹렬한 궤적과, 결코 평탄해질 수 없는 인간의 '비대칭적 번식 구조'가 정면으로 충돌하며 뿜어내는 파괴적인 물리학적 파열음이다.

가장 고도로 발전하고 가장 공정해진 선진국일수록 아이의 울음소리가 멈추는 이유는 그들이 부유해서가 아니다. 그들이 우주 역사상 가장 맹렬하고 완벽하게 닦인 '대칭적 단일 경쟁 트랙'을 완성해 냈기 때문이다. 문명이 비대칭성을 매끄럽게 깎아내려 할수록, 번식이라는 엇갈린 디스턴스는 그 매끄러운 표면 위에서 산산조각 날 수밖에 없는 웅장한 비극적 결말을 맞이하고 있는 것이다.
